\chapter{正式估值、Street与模型审计}

第五章不是从39只中简单选出“前三名”，而是展示正式估值审计集，并把其余标的的降级原因完整公开。39只优先池先经过当前价格、概率空间、证据质量、研报日期、外部目标价和原始PDF路径六道门槛，再分成正式模型、条件性高空间观察、条件性安全边际观察和估值风险观察。

选择桥结果：条件性高空间观察 12只、正式审计模型 3只、条件性安全边际观察 9只、估值风险/观察 15只。正式审计集包含两只正向空间模型和一只证据链完整但概率目标低于现价的下行风险对照；这三只不代表39只中只有三只值得研究。

\begin{exhibitbox}[Exhibit：三只正式当前价模型]
\scriptsize
\begin{tabularx}{\textwidth}{L{1.0cm}L{1.6cm}L{1.5cm}R{0.85cm}R{0.85cm}R{0.85cm}R{0.85cm}R{0.85cm}X}
\toprule
\textbf{代码} & \textbf{公司} & \textbf{方法} & \textbf{熊} & \textbf{基准} & \textbf{牛} & \textbf{概率值} & \textbf{空间} & \textbf{动作} \\
\midrule
000703 & 恒逸石化 & cycle-adjusted deducted-EPS PE & 10.86 & 25.59 & 40.41 & 23.42 & 61.7\% & 高空间验证 \\
002379 & 宏桥控股 & cycle-adjusted deducted-EPS PE & 14.00 & 26.24 & 40.65 & 25.80 & 38.3\% & 高空间验证 \\
000155 & 川能动力 & deducted-EPS scenario PE & 5.43 & 9.03 & 13.67 & 10.04 & -9.5\% & 估值观察 \\
\bottomrule
\end{tabularx}
\normalsize
\end{exhibitbox}

\section{逐票审计}
\subsection{恒逸石化（000703）}
现价14.48元，H1归母/扣非中值57.5/57.2亿元；Q1经营现金流-25.7亿元。模型熊/基准/牛10.86/25.59/40.41元，概率值23.42元，空间61.7\%。最新研报状态current，目标价字段17.05；动作为高空间验证。

\subsection{宏桥控股（002379）}
现价18.66元，H1归母/扣非中值155.0/154.0亿元；Q1经营现金流92.4亿元。模型熊/基准/牛14.00/26.24/40.65元，概率值25.80元，空间38.3\%。最新研报状态current，目标价字段29.00；动作为高空间验证。

\subsection{川能动力（000155）}
现价11.09元，H1归母/扣非中值6.8/6.8亿元；Q1经营现金流-3.2亿元。模型熊/基准/牛5.43/9.03/13.67元，概率值10.04元，空间-9.5\%。最新研报状态current，目标价字段20.43；动作为估值观察。

\section{概率目标的计算公式}

\begin{exhibitbox}[概率目标计算公式与参数顺序]
\small
\begin{itemize}
\item 股本（亿股） = 总市值（亿元） / 当前股价。
\item H1扣非每股收益 = H1扣非归母净利润中值（亿元） / 股本。
\item 全年情景每股收益 = H1扣非每股收益 ×（1 + H2/H1利润转换比例）。
\item 情景估值 = 情景每股收益 × 对应行业估值倍数。金融股若有Q1每股净资产，则使用 65\% PB估值 + 35\% PE估值的混合结果。
\item 熊值下限 = min（原始熊值，当前股价 × 75\%），确保模型保留真实下行空间。
\item House概率价值 = 熊值 × 30\% + 基准值 × 50\% + 牛值 × 20\%。
\item 最终概率目标 = House概率价值 ×（1 - 外部锚权重） + 外部目标价 × 外部锚权重。
\item 上涨/下跌空间 = 最终概率目标 / 当前股价 - 1。当前模型的市场情绪锚权重为 0\%。
\end{itemize}
\normalsize
\end{exhibitbox}

参数来源边界：H1扣非利润来自公司业绩预告；股本由当前总市值与当前股价反推；H2/H1比例按行业类别设定，并在启动确认或资金观察阶段对基准、牛市比例上调；行业PE区间来自预设的业务类型匹配区间。该模型是筛选和验证模型，不是公司管理层指引。

\section{三只正式模型的代入计算}

\subsection{恒逸石化（000703）}
股本 = 553.36亿元 / 14.48元 = 38.2155亿股；H1扣非EPS = 57.20亿元 / 38.2155亿股 = 1.4968元。
H2/H1转换比例为熊/基准/牛 0.55/0.90/1.25；情景EPS为 2.3200/2.8439/3.3677元；对应原始情景值为 13.92/25.59/40.41元。
熊值经过当前价下限约束后为 10.86元；House概率价值 = 10.86 × 0.30 + 25.59 × 0.50 + 40.41 × 0.20 = 24.13元。
外部目标价为 17.05元，外部锚权重为 10\%；最终概率目标 = 24.13 × 90\% + 17.05 × 10\% = 23.42元；空间 = 23.42 / 14.48 - 1 = 61.7\%。

\subsection{宏桥控股（002379）}
股本 = 2431.61亿元 / 18.66元 = 130.3114亿股；H1扣非EPS = 154.00亿元 / 130.3114亿股 = 1.1818元。
H2/H1转换比例为熊/基准/牛 0.55/0.85/1.15；情景EPS为 1.8318/2.1863/2.5408元；对应原始情景值为 14.65/26.24/40.65元。
熊值经过当前价下限约束后为 14.00元；House概率价值 = 14.00 × 0.30 + 26.24 × 0.50 + 40.65 × 0.20 = 25.45元。
外部目标价为 29.00元，外部锚权重为 10\%；最终概率目标 = 25.45 × 90\% + 29.00 × 10\% = 25.80元；空间 = 25.80 / 18.66 - 1 = 38.3\%。

\subsection{川能动力（000155）}
股本 = 204.74亿元 / 11.09元 = 18.4617亿股；H1扣非EPS = 6.75亿元 / 18.4617亿股 = 0.3656元。
H2/H1转换比例为熊/基准/牛 0.65/0.90/1.20；情景EPS为 0.6033/0.6947/0.8044元；对应原始情景值为 5.43/9.03/13.67元。
熊值经过当前价下限约束后为 5.43元；House概率价值 = 5.43 × 0.30 + 9.03 × 0.50 + 13.67 × 0.20 = 8.88元。
外部目标价为 20.43元，外部锚权重为 10\%；最终概率目标 = 8.88 × 90\% + 20.43 × 10\% = 10.04元；空间 = 10.04 / 11.09 - 1 = -9.5\%。

\section{模型边界}

\section{39只优先池到正式模型的桥接}

\scriptsize
\begin{longtable}{L{1.15cm}L{1.55cm}L{2.15cm}R{1.15cm}R{1.05cm}L{3.0cm}}
\toprule
\textbf{代码} & \textbf{标的} & \textbf{层级} & \textbf{空间} & \textbf{外部锚} & \textbf{未升级原因/使用方式} \\
\midrule
\endfirsthead
\toprule
\textbf{代码} & \textbf{标的} & \textbf{层级} & \textbf{空间} & \textbf{外部锚} & \textbf{未升级原因/使用方式} \\
\midrule
\endhead
000623 & 吉林敖东 & 条件性高空间观察 & 317.0\% & House替代估值 & 证据质量不足以支持正式模型；最新研报缺失或早于2026-04-01；没有可审计的外部目标价/合理价值区间；没有对应的本地原始研报证据路径 \\
600739 & 辽宁成大 & 条件性高空间观察 & 290.4\% & 历史/弱来源目标22.79元（零权） & 证据质量不足以支持正式模型；最新研报缺失或早于2026-04-01 \\
600150 & 中国船舶 & 条件性高空间观察 & 142.9\% & House替代估值 & 没有可审计的外部目标价/合理价值区间；没有对应的本地原始研报证据路径 \\
601600 & 中国铝业 & 条件性高空间观察 & 62.8\% & House替代估值 & 没有可审计的外部目标价/合理价值区间；没有对应的本地原始研报证据路径 \\
000703 & 恒逸石化 & 正式审计模型 & 61.7\% & 历史/弱来源目标17.05元（零权） & 正式模型：当前价格、情景估值、外部锚和原始研报证据均可复核 \\
601061 & 中信金属 & 条件性高空间观察 & 57.7\% & 历史/弱来源目标11.25元（零权） & 证据质量不足以支持正式模型；最新研报缺失或早于2026-04-01 \\
601360 & 三六零 & 条件性高空间观察 & 45.6\% & 历史/弱来源目标14.44元（零权） & 证据质量不足以支持正式模型；最新研报缺失或早于2026-04-01 \\
002379 & 宏桥控股 & 正式审计模型 & 38.3\% & 历史/弱来源目标29.00元（零权） & 正式模型：当前价格、情景估值、外部锚和原始研报证据均可复核 \\
600026 & 中远海能 & 条件性高空间观察 & 34.8\% & House替代估值 & 没有可审计的外部目标价/合理价值区间；没有对应的本地原始研报证据路径 \\
002048 & 宁波华翔 & 条件性高空间观察 & 33.0\% & House替代估值 & 证据质量不足以支持正式模型；最新研报缺失或早于2026-04-01；没有可审计的外部目标价/合理价值区间；没有对应的本地原始研报证据路径 \\
000737 & 北方铜业 & 条件性高空间观察 & 20.8\% & House替代估值 & 证据质量不足以支持正式模型；最新研报缺失或早于2026-04-01；没有可审计的外部目标价/合理价值区间；没有对应的本地原始研报证据路径 \\
000567 & 海德股份 & 条件性高空间观察 & 20.8\% & House替代估值 & 证据质量不足以支持正式模型；最新研报缺失或早于2026-04-01；没有可审计的外部目标价/合理价值区间；没有对应的本地原始研报证据路径 \\
601336 & 新华保险 & 条件性高空间观察 & 19.7\% & House替代估值 & 没有可审计的外部目标价/合理价值区间；没有对应的本地原始研报证据路径 \\
002602 & 世纪华通 & 条件性高空间观察 & 15.8\% & House替代估值 & 没有可审计的外部目标价/合理价值区间；没有对应的本地原始研报证据路径 \\
002558 & 巨人网络 & 条件性安全边际观察 & 15.0\% & House替代估值 & 没有可审计的外部目标价/合理价值区间；没有对应的本地原始研报证据路径 \\
300014 & 亿纬锂能 & 条件性安全边际观察 & 12.1\% & House替代估值 & 没有可审计的外部目标价/合理价值区间；没有对应的本地原始研报证据路径 \\
601688 & 华泰证券 & 条件性安全边际观察 & 9.4\% & House替代估值 & 没有可审计的外部目标价/合理价值区间；没有对应的本地原始研报证据路径 \\
000042 & 中洲控股 & 条件性安全边际观察 & 6.5\% & 历史/弱来源目标11.84元（零权） & 证据质量不足以支持正式模型；最新研报缺失或早于2026-04-01 \\
600489 & 中金黄金 & 条件性安全边际观察 & 5.2\% & House替代估值 & 没有可审计的外部目标价/合理价值区间；没有对应的本地原始研报证据路径 \\
000975 & 山金国际 & 条件性安全边际观察 & 5.0\% & House替代估值 & 没有可审计的外部目标价/合理价值区间；没有对应的本地原始研报证据路径 \\
002460 & 赣锋锂业 & 条件性安全边际观察 & 4.2\% & House替代估值 & 没有可审计的外部目标价/合理价值区间；没有对应的本地原始研报证据路径 \\
000987 & 越秀资本 & 条件性安全边际观察 & 2.5\% & 历史/弱来源目标9.44元（零权） & 证据质量不足以支持正式模型；最新研报缺失或早于2026-04-01 \\
600120 & 浙江东方 & 条件性安全边际观察 & 2.2\% & House替代估值 & 证据质量不足以支持正式模型；最新研报缺失或早于2026-04-01；没有可审计的外部目标价/合理价值区间；没有对应的本地原始研报证据路径 \\
601377 & 兴业证券 & 估值风险/观察 & -0.8\% & House替代估值 & 没有可审计的外部目标价/合理价值区间；没有对应的本地原始研报证据路径；概率目标低于当前价，仅保留为下行风险对照 \\
000858 & 五粮液 & 估值风险/观察 & -1.5\% & House替代估值 & 没有可审计的外部目标价/合理价值区间；没有对应的本地原始研报证据路径；概率目标低于当前价，仅保留为下行风险对照 \\
000301 & 东方盛虹 & 估值风险/观察 & -3.7\% & House替代估值 & 没有可审计的外部目标价/合理价值区间；没有对应的本地原始研报证据路径；概率目标低于当前价，仅保留为下行风险对照 \\
002756 & 永兴材料 & 估值风险/观察 & -6.2\% & House替代估值 & 没有可审计的外部目标价/合理价值区间；没有对应的本地原始研报证据路径；概率目标低于当前价，仅保留为下行风险对照 \\
600918 & 中泰证券 & 估值风险/观察 & -6.3\% & House替代估值 & 证据质量不足以支持正式模型；最新研报缺失或早于2026-04-01；没有可审计的外部目标价/合理价值区间；没有对应的本地原始研报证据路径；概率目标低于当前价，仅保留为下行风险对照 \\
002466 & 天齐锂业 & 估值风险/观察 & -6.4\% & House替代估值 & 没有可审计的外部目标价/合理价值区间；没有对应的本地原始研报证据路径；概率目标低于当前价，仅保留为下行风险对照 \\
600346 & 恒力石化 & 估值风险/观察 & -6.8\% & House替代估值 & 没有可审计的外部目标价/合理价值区间；没有对应的本地原始研报证据路径；概率目标低于当前价，仅保留为下行风险对照 \\
600688 & 上海石化 & 估值风险/观察 & -8.1\% & 历史/弱来源目标4.06元（零权） & 证据质量不足以支持正式模型；最新研报缺失或早于2026-04-01；概率目标低于当前价，仅保留为下行风险对照 \\
000155 & 川能动力 & 正式审计模型 & -9.5\% & 历史/弱来源目标20.43元（零权） & 正式下行风险对照：证据链完整，但概率目标低于当前价 \\
600037 & 歌华有线 & 估值风险/观察 & -17.3\% & House替代估值 & 证据质量不足以支持正式模型；最新研报缺失或早于2026-04-01；没有可审计的外部目标价/合理价值区间；没有对应的本地原始研报证据路径；概率目标低于当前价，仅保留为下行风险对照 \\
002738 & 中矿资源 & 估值风险/观察 & -20.4\% & House替代估值 & 没有可审计的外部目标价/合理价值区间；没有对应的本地原始研报证据路径；概率目标低于当前价，仅保留为下行风险对照 \\
000426 & 兴业银锡 & 估值风险/观察 & -22.9\% & House替代估值 & 没有可审计的外部目标价/合理价值区间；没有对应的本地原始研报证据路径；概率目标低于当前价，仅保留为下行风险对照 \\
001339 & 智微智能 & 估值风险/观察 & -26.7\% & House替代估值 & 没有可审计的外部目标价/合理价值区间；没有对应的本地原始研报证据路径；概率目标低于当前价，仅保留为下行风险对照 \\
601069 & 西部黄金 & 估值风险/观察 & -44.5\% & House替代估值 & 证据质量不足以支持正式模型；最新研报缺失或早于2026-04-01；没有可审计的外部目标价/合理价值区间；没有对应的本地原始研报证据路径；概率目标低于当前价，仅保留为下行风险对照 \\
600095 & 湘财股份 & 估值风险/观察 & -46.2\% & House替代估值 & 证据质量不足以支持正式模型；最新研报缺失或早于2026-04-01；没有可审计的外部目标价/合理价值区间；没有对应的本地原始研报证据路径；概率目标低于当前价，仅保留为下行风险对照 \\
600292 & 电投水电 & 估值风险/观察 & -47.5\% & House替代估值 & 证据质量不足以支持正式模型；最新研报缺失或早于2026-04-01；没有可审计的外部目标价/合理价值区间；没有对应的本地原始研报证据路径；概率目标低于当前价，仅保留为下行风险对照 \\
\bottomrule
\end{longtable}
\normalsize

正式模型的选择重点是“证据是否足以支撑当前价格下的可审计估值”，不是单纯追求概率空间最大。中洲控股、吉林敖东、辽宁成大和中国船舶等标的模型空间很高，但分别受到外部锚缺失、研报过旧、证据质量不足或原始目标价缺失的约束，因此进入条件性观察而非正式Street锚模型。川能动力虽然证据链完整，但概率目标低于现价，故作为下行风险对照，不作为正向推荐。

本轮动态模型是全量筛选后的可复算模型层，不等于对每家公司都完成了深度公司研究。正式模型仍需独立IC复核；条件性观察池的升级条件已经在第四章逐票写明。所有陈旧、失败或弱证据研报不作为正向Street锚。
\section{估值恢复模型：三六零与中洲控股}

这两只标的原先的H1扭亏/项目结算PE外推会产生不可解释的极端空间，因此本轮改用估值恢复模型。恢复模型不把弱来源直接当作Street目标，而是将公开机构共识、可比公司方法、市场隐含估值和替代估值方法分层使用。
\subsection{三六零（601360）}
恢复状态：recovered\_conditional；主方法：PS on 2026E revenue；可比公司：昆仑万维、启明星辰、金山办公；置信度：conditional\_medium。
情景区间为10.28/13.70/15.07元；外部锚说明：公开共识均值14.44元，赋予10\%权重；公开共识来源质量为auditable\_consensus\_snapshot。
2026E revenue 95.91亿元 / 70.00亿股 × PS multiple
概率价值 = 10.28×0.30 + 13.70×0.50 + 15.07×0.20 = 12.95元。
最终条件性目标 = 12.95×0.90 + 14.44×0.10 = 13.10元，空间45.6\%。
市场隐含指标：AI+安全重估，但资金参与度偏弱、融资余额低于一年30\%分位。恢复模型置信度为conditional\_medium，不等同于正式Street目标。
正常化分母：2026E revenue consensus and recurring AI/security monetization, not H1 turnaround EPS。升级条件：2026E revenue near 95.91亿元, recurring AI/security monetization, EPS consensus >=0.06, cash flow turns positive。降级条件：AI product monetization fails, 2026E EPS falls below 0.03, OCF remains negative, or PS falls below peer-supported range。
\subsection{中洲控股（000042）}
恢复状态：recovered\_conditional；主方法：normalized EPS / PE with project settlement scenarios；可比公司：保利发展、招商蛇口、滨江集团、越秀地产；置信度：conditional\_low\_medium。
情景区间为4.80/8.80/13.50元；外部锚说明：弱来源/历史目标11.84元，明确赋予0\%权重；公开共识来源质量为auditable\_public\_snapshot。
bear 0.80×6=4.80; base 1.10×8=8.80; bull 1.50×9=13.50; H1 1.4891元 EPS is not annualized directly
概率价值 = 4.80×0.30 + 8.80×0.50 + 13.50×0.20 = 8.54元。
无正权重外部目标锚，最终条件性目标等于概率价值8.54元，空间6.5\%。
市场隐含指标：业绩预告刺激有限，7/15主力净流出约1473万元、连续2日减仓；地产高负债和项目结算风险压制估值。恢复模型置信度为conditional\_low\_medium，不等同于正式Street目标。
正常化分母：recurring/normalized EPS after project settlement and tax/interest normalization。升级条件：2026E EPS at least 1.10, project settlement/tax provisions confirmed, debt and cash-flow pressure improves, and price holds above 8。降级条件：settlement profit reverses, land appreciation tax/interest provisions rise, sales/collection weaken, or net debt/liquidity worsens。
\subsection{赣锋锂业（002460）}
恢复状态：recovered\_conditional；主方法：cycle-adjusted EPS / PE；可比公司：天齐锂业、中矿资源、盛新锂能；置信度：conditional\_medium。
情景区间为40.00/53.40/74.90元；外部锚说明：没有可审计正权外部目标，采用House替代估值；公开共识来源质量为未形成公开目标价共识。
2026E EPS 4.45元来自当期原始研报；熊/基准/牛采用锂价持续性和资源兑现敏感的10/12/14倍PE
概率价值 = 40.00×0.30 + 53.40×0.50 + 74.90×0.20 = 53.68元。
无正权重外部目标锚，最终条件性目标等于概率价值53.68元，空间4.2\%。
市场隐含指标：锂价上行带来盈利弹性，但周期持续性和海外资源兑现仍需验证。恢复模型置信度为conditional\_medium，不等同于正式Street目标。
正常化分母：2026E cycle-adjusted EPS, not H1 turnaround annualization。升级条件：锂价、产量、资源自给率和H2经营现金流共同验证2026E EPS 4.45元。降级条件：锂价回落、资源项目延期、库存减值或现金流弱于利润。
\subsection{歌华有线（600037）}
恢复状态：recovered\_conditional；主方法：normalized EPS / PE with turnaround discount；可比公司：东方明珠、电广传媒、新媒股份；置信度：conditional\_low。
情景区间为2.40/5.40/11.00元；外部锚说明：没有可审计正权外部目标，采用House替代估值；公开共识来源质量为未形成公开目标价共识。
公开预测EPS 0.03/0.03/0.05元，按扭亏持续性和高估值风险采用120/180/220倍PE
概率价值 = 2.40×0.30 + 5.40×0.50 + 11.00×0.20 = 5.62元。
无正权重外部目标锚，最终条件性目标等于概率价值5.62元，空间-17.3\%。
市场隐含指标：扭亏预期已被部分交易，通信业务转型和盈利持续性仍弱。恢复模型置信度为conditional\_low，不等同于正式Street目标。
正常化分母：2026E recurring EPS after turnaround, not H1 profit annualization。升级条件：通信业务收入、现金流和全年EPS达到公开预测并连续两个季度验证。降级条件：扭亏由一次性收益驱动、通信业务不及预期或现金流转弱。
\subsection{宁波华翔（002048）}
恢复状态：recovered\_conditional；主方法：normalized EPS / PE；可比公司：华域汽车、福耀玻璃、拓普集团；置信度：conditional\_medium。
情景区间为19.20/29.64/39.60元；外部锚说明：没有可审计正权外部目标，采用House替代估值；公开共识来源质量为未形成公开目标价共识。
2026E EPS 1.976元来自研报，按汽车零部件业务和机器人弹性采用12/15/18倍PE
概率价值 = 19.20×0.30 + 29.64×0.50 + 39.60×0.20 = 28.50元。
无正权重外部目标锚，最终条件性目标等于概率价值28.50元，空间33.0\%。
市场隐含指标：汽车零部件和机器人业务提供弹性，但预测时点较旧。恢复模型置信度为conditional\_medium，不等同于正式Street目标。
正常化分母：2026E normalized EPS after auto-cycle and robot-order verification。升级条件：汽车订单、机器人批量供货和H2扣非利润验证2026E EPS 1.976元。降级条件：主机厂降价、机器人订单延迟或扣非利润低于预期。
\subsection{上海石化（600688）}
恢复状态：recovered\_conditional；主方法：PB / cycle-adjusted asset value；可比公司：中国石化、恒力石化、华锦股份；置信度：conditional\_low\_medium。
情景区间为2.01/2.68/3.35元；外部锚说明：弱来源/历史目标4.06元，明确赋予0\%权重；公开共识来源质量为未形成公开目标价共识。
Q1 BPS 2.2324元 × 0.90/1.20/1.50倍PB；历史4.06元目标权重为0\%
概率价值 = 2.01×0.30 + 2.68×0.50 + 3.35×0.20 = 2.61元。
无正权重外部目标锚，最终条件性目标等于概率价值2.61元，空间-8.1\%。
市场隐含指标：炼化周期和现金流承压，历史目标价不能作为当前Street锚。恢复模型置信度为conditional\_low\_medium，不等同于正式Street目标。
正常化分母：Q1 BPS and cycle-adjusted asset value。升级条件：炼化价差、经营现金流和资产负债表改善，且新一期机构预测恢复。降级条件：炼化价差继续收窄、现金流恶化或资产减值上升。
\subsection{德明利（001309）}
恢复状态：recovered\_conditional；主方法：forward EPS / PE with memory-cycle discount；可比公司：江波龙、兆易创新、北京君正；置信度：conditional\_medium。
情景区间为330.00/513.00/630.00元；外部锚说明：没有可审计正权外部目标，采用House替代估值；公开共识来源质量为未形成公开目标价共识。
2026E EPS 39.16元来自原始研报；按存储周期和估值消化采用11/13.1/14倍PE
概率价值 = 330.00×0.30 + 513.00×0.50 + 630.00×0.20 = 481.50元。
无正权重外部目标锚，最终条件性目标等于概率价值481.50元，空间-27.3\%。
市场隐含指标：存储价格和企业级产品景气强，但周期反转后利润回撤风险高。恢复模型置信度为conditional\_medium，不等同于正式Street目标。
正常化分母：2026E memory-cycle-adjusted EPS, not H1 annualized EPS。升级条件：存储价格、企业级产品出货和H2现金流确认2026E EPS 39.16元。降级条件：存储价格反转、库存减值、客户集中或经营现金流持续为负。
\subsection{天华新能（300390）}
恢复状态：recovered\_conditional；主方法：cycle-adjusted EPS / PE；可比公司：赣锋锂业、盛新锂能、天齐锂业；置信度：conditional\_medium。
情景区间为56.40/79.80/102.40元；外部锚说明：没有可审计正权外部目标，采用House替代估值；公开共识来源质量为未形成公开目标价共识。
2026E EPS约5.70元来自原始研报；可比公司PE约16/13/11倍，采用12/14/16倍情景
概率价值 = 56.40×0.30 + 79.80×0.50 + 102.40×0.20 = 77.30元。
无正权重外部目标锚，最终条件性目标等于概率价值77.30元，空间28.8\%。
市场隐含指标：锂盐价格恢复和大客户绑定带来弹性，但周期延续性仍需验证。恢复模型置信度为conditional\_medium，不等同于正式Street目标。
正常化分母：2026E cycle-adjusted EPS and lithium-price sensitivity。升级条件：氢氧化锂价格、客户采购、产量和H2经营现金流确认2026E EPS 5.70元。降级条件：锂价回落、客户集中风险暴露或经营现金流继续弱于利润。
\subsection{恩捷股份（002812）}
恢复状态：recovered\_conditional；主方法：forward EPS / PE with capacity-cycle normalization；可比公司：星源材质、中材科技、沧州明珠；置信度：conditional\_medium。
情景区间为21.60/103.20/126.20元；外部锚说明：原始研报目标103.00元，赋予10\%权重；公开共识来源质量为未形成公开目标价共识。
熊情景用2026E EPS 1.80元；基准/牛情景用2027E EPS 5.16/6.31元×20倍，目标价103元作为原始PDF交叉锚
概率价值 = 21.60×0.30 + 103.20×0.50 + 126.20×0.20 = 83.32元。
无正权重外部目标锚，最终条件性目标等于概率价值85.29元，空间55.2\%。
市场隐含指标：隔膜供需改善和涨价预期强化，但2027E利润跳升尚需订单和价格验证。恢复模型置信度为conditional\_medium，不等同于正式Street目标。
正常化分母：2027E normalized EPS after separator price and utilization recovery。升级条件：隔膜涨价、产能利用率、海外/3C订单和H2现金流确认2027E EPS 5.16元。降级条件：供需修复不及预期、价格回落、客户集中或产能利用率不足。
\subsection{盛新锂能（002240）}
恢复状态：recovered\_conditional；主方法：cycle-adjusted EPS / PE；可比公司：赣锋锂业、天齐锂业、天华新能；置信度：conditional\_medium。
情景区间为21.60/39.20/59.04元；外部锚说明：没有可审计正权外部目标，采用House替代估值；公开共识来源质量为未形成公开目标价共识。
2026E EPS 2.45元来自原始研报，按锂价和自有矿兑现采用12/16/18倍PE
概率价值 = 21.60×0.30 + 39.20×0.50 + 59.04×0.20 = 37.89元。
无正权重外部目标锚，最终条件性目标等于概率价值37.89元，空间21.8\%。
市场隐含指标：锂价和自有矿贡献改善，但H1经营现金流为负、周期持续性仍是核心变量。恢复模型置信度为conditional\_medium，不等同于正式Street目标。
正常化分母：2026E cycle-adjusted EPS with lithium-price and cash-flow bridge。升级条件：锂价、出货、自有矿贡献和H2经营现金流验证2026E EPS 2.45元。降级条件：锂价回落、现金流持续为负或资源项目贡献不及预期。
\subsection{苏州高新（600736）}
恢复状态：recovered\_conditional；主方法：normalized EPS / PE with PB cross-check；可比公司：张江高科、陆家嘴、外高桥；置信度：conditional\_low\_medium。
情景区间为4.20/7.70/10.40元；外部锚说明：没有可审计正权外部目标，采用House替代估值；公开共识来源质量为未形成公开目标价共识。
2026E EPS 0.14元来自原始研报，按地产结算和产业园转型采用35/55/65倍PE
概率价值 = 4.20×0.30 + 7.70×0.50 + 10.40×0.20 = 7.19元。
无正权重外部目标锚，最终条件性目标等于概率价值7.19元，空间4.7\%。
市场隐含指标：产业园和投资收益提供支撑，但地产结算波动和现金流承压。恢复模型置信度为conditional\_low\_medium，不等同于正式Street目标。
正常化分母：2026E normalized EPS after property settlement and investment-income normalization。升级条件：产业园利润、地产结算、投资收益和经营现金流共同验证2026E EPS 0.14元。降级条件：地产结算延期、投资收益回落或经营现金流继续为负。
\subsection{渤海租赁（000415）}
恢复状态：recovered\_conditional；主方法：normalized EPS / PE with PB cross-check；可比公司：中银航空租赁、中国飞机租赁、国银租赁；置信度：conditional\_medium。
情景区间为3.00/4.62/5.95元；外部锚说明：没有可审计正权外部目标，采用House替代估值；公开共识来源质量为未形成公开目标价共识。
2026E EPS 0.77元来自原始研报，按航空租赁资产质量和杠杆采用5/6/7倍PE
概率价值 = 3.00×0.30 + 4.62×0.50 + 5.95×0.20 = 4.40元。
无正权重外部目标锚，最终条件性目标等于概率价值4.40元，空间-1.1\%。
市场隐含指标：航空租赁景气改善，但资产质量、汇率和杠杆仍压制估值。恢复模型置信度为conditional\_medium，不等同于正式Street目标。
正常化分母：2026E normalized EPS after lease income, asset quality and FX normalization。升级条件：租赁收入、资产质量、汇率和H2现金流确认2026E EPS 0.77元。降级条件：飞机减值、汇率损失、融资成本或现金流风险上升。
